\documentclass[12pt,a4paper]{article}  % Déclare la classe du document.

%  ------------------------------------------  Modules usuels  --------------------------------------
%  --------------------------------------------------------------------------------------------------------

\usepackage{amsmath} 	
\usepackage{amssymb} 
\usepackage{listings}         
\usepackage{numprint} 
\usepackage{pifont}

%  --------------------------------------------------------------------------------------------------------
% ----------------------  packages pour fontes et caractères spéciaux  ---------------------
%  --------------------------------------------------------------------------------------------------------

\usepackage{eurosym}   
\usepackage{amsfonts}   
\usepackage{amstext}       
\usepackage{amsfonts}	
\usepackage{mathrsfs}	 

\usepackage{pstricks,pst-all,pstricks-add,pst-tree,pst-node,pst-3dplot}

%------------------------------------------------------------------------------------------------------
%  --------------------------  Francisation et gestion des caractères sous mac --------------------------------------------------------------------------------------------------------

\usepackage[francais]{babel}   
\usepackage[utf8]{inputenc}
\usepackage[OT1,T1]{fontenc} 
%  --------------------------------------------------------------------------------------------------------
% ------------------------------------Packages  de  mise en page ---------------------------
%--------------------------------------------------------------------------------------------------------

\usepackage{fancyhdr}	% en-têtes et pieds de pages
\usepackage{amsfonts}

 \usepackage{multicol} % permet l'utilisation du multicolonnage dans le texte. La syntaxe est \begin{multicols}{nb de colonnes}    ......  \end{multicols}
 
\usepackage{calc}            % permet de faire un minimum de calcul sur les arguments des commandes et macros (utilisé dans la def de l'environnement preuve )
\usepackage{epic}            % extension et simplification de  l'environnement picture
\usepackage{graphicx}    % permet d'inclure des graphique externe dans un document
\usepackage{variations}   % permet d'obtenir de joli tableau de variations
\usepackage{enumerate}  %permet de faire des listes personnalisées
\usepackage{arydshln}     %permet d'avoir des lignes en pointillés
\usepackage{framed}  
\usepackage{multido} 


%  --------------------------------------  MACROS USUELLES  ----------------------------------------

\newcommand{\R}{\mathbb{R}} 
\newcommand{\N}{\mathbb{N}} 
\newcommand{\D}{\mathbb{D}} 
\newcommand{\Z}{\mathbb{Z}} 
\newcommand{\Q}{\mathbb{Q}} 
\newcommand{\C}{\mathbb{C}} 

%  -------------------------------------  MACROS SPECIFIQUES ------------------------------------


\newcommand{\raisedrule}[2][0em]{\leaders\hbox{\rule[#1]{1pt}{#2}}\hfill}


\newcounter{numexo}

\setcounter{numexo}{1}

\newcommand{\exercice}[1]{\vskip 0.6cm\noindent\rule[+2pt]{1cm}{0.5pt}\ {\large\bf \ Exercice \thenumexo {  #1 }}\raisedrule[+2pt]{0.5pt}\null\addtocounter{numexo}{1}\par\vskip 0,4cm}


\newcommand{\lignes}[1]{
\setlength{\unitlength}{1mm}
\vskip 0,4cm\noindent
\multido{\i=0+1}{#1}{%
     \noindent
     \begin{picture}(150,5)
     \linethickness{0,005mm}
     \put(-5,0){\dashline{1}(175,5)(5,5)}
     \end{picture}\vskip 0,001cm
     }
\vskip 0cm
}




% ----------------------------------------- FORMAT DE PAGE  -----------------------------------------

\usepackage[left = 1.3cm,
	 	      right = 1.3cm,
		      top = 1.7cm,
		      headsep = 8mm,
		      bottom = 1.7cm,
		      headheight = 14.2 pt
		      ]{geometry}

\setlength{\parindent}{0mm}


\everymath{\displaystyle} 
\renewcommand{\thesection}{\Roman{section}} 
\author{}
\title{}
\date{}
\pagestyle{fancy}

 
\lhead{1ère - spécialité maths} 


\rhead{\textcolor{red}{Correction - Rattrapage - DS n$^{\circ}$5 - 23/02/21}}		


\fancyfoot[C]{} 
\fancyfoot[L]{} 
\fancyfoot[R]{} 

\lstset{language=Python}  

\lstset{frame = single
framerule = 1pt
showspaces = false,
showstringspaces = false,
numbers=left, numberstyle=\tiny}

% ----------------------------------------- correction  -----------------------------------------
\usepackage[most,skins,breakable]{tcolorbox}
\usepackage{enumitem}

\newcommand{\corr}[1]{\fcolorbox{red}{yellow!30}{#1}}

\newcommand{\bareme}[1]{\tikz\node[circle,draw,inner sep=1pt,outer sep=1pt, color=red]{\textbf{#1}};}

\newtcolorbox{corrBox}{enhanced,boxrule=1pt,nobeforeafter,boxsep=0pt,left=\fboxsep,right=\fboxsep,colback=yellow!30,colframe=red!75!black}

\newlist{enumMet}{enumerate}{2}
\setlist[enumMet,1]{label=\arabic*), font=\bf\color{blue!75!black}}
\setlist[enumMet,2]{label=\alph*), font=\bf\color{blue!75!black}}


\begin{document}
\graphicspath{{images/}}

\begin{center}

\fbox{\LARGE{\rule{0.5cm}{0pt}\rule[-0.25cm]{0pt}{0.9cm} \textcolor{red}{Correction - Rattrapage DS N$^{\circ}$5}	\rule{0.5cm}{0pt}}}
\end{center}


\begin{center}\textbf{1h - Calculatrice autorisée - Barème indicatif}\end{center}

%\begin{center}\textbf{Les élèves avec un tier-temps ne traitent pas les questions avec  le symbole \ding{46}}\end{center}

\begin{center}
%On veillera à ENCADRER ses résultats et soigner sa copie.\\
Les résultats doivent être justifiés par des calculs (au moins 1 étape intermédiaire).
\end{center}
\vspace{.05cm} 

\vspace{-0.3cm}
\exercice{ - définition du produit scalaire - ( 3 points )}
On considère 3 points $A, B$ et $C .$ Dans chacun des cas suivants, calculer $\overrightarrow{A B} \cdot \overrightarrow{A C}$ ( on donnera la valeur exacte).
\begin{multicols}{2}
\begin{enumMet}
    \item $A B=4, A C=5$ et $\widehat{A B C}=\dfrac{\pi}{6}$ \quad \bareme{1,5}\\
    \begin{corrBox}
        $\overrightarrow{A B} \cdot \overrightarrow{A C}=20 \times \cos \left(\frac{\pi}{6}\right)=20 \times \frac{\sqrt{3}}{2}=10 \sqrt{3}$
    \end{corrBox}
    
    \columnbreak
    \item $A B=5, A C=7$ et $\widehat{A B C}=\dfrac{3\pi}{4}$ \quad \bareme{1,5}\\
    \begin{corrBox}
        $\overrightarrow{A B} \cdot \overrightarrow{A C}=7 \times 5 \times \cos \left(\frac{3 \pi}{4}\right)=-35 \times \frac{\sqrt{2}}{2}=-\frac{35 \sqrt{2}}{2}$
    \end{corrBox}
\end{enumMet}
\end{multicols}

\exercice{ - produit scalaire dans un triangle - ( 2 points )}

Soit un triangle $\mathrm{ABC}$ de longueur $\mathrm{AB}=4, \mathrm{AC}=6$ et $\mathrm{BC}=9$\\
Calculer le produit scalaire $\overrightarrow{B A} \cdot \overrightarrow{B C}$

\begin{corrBox}
    $$\overrightarrow{B A} \cdot \overrightarrow{B C}=\frac{1}{2}\left(B A^{2}+B C^{2}-A C^{2}\right)\quad \bareme{1}\quad=\frac{1}{2}\left(4^{2}+9^{2}-6^{2}\right)=\frac{1}{2}(16+81-36)$$
$$\overrightarrow{B A} \cdot \overrightarrow{B C}=\frac{61}{2}=30,5\quad \bareme{1}$$
\end{corrBox}

\exercice{ - produit scalaire dans un repère - ( 4 points )}

Soit les points $A(3 ; 10), B(5 ; 5)$ et $C(0 ; 3)$
\begin{enumMet}
    \item Calculer le produit scalaire $\overrightarrow{B A} \cdot \overrightarrow{B C}$\\
    \begin{corrBox}
        $\overrightarrow{B A} \begin{pmatrix} 3-5\\ 10-5 \end{pmatrix}=\overrightarrow{B A} \begin{pmatrix} -2\\ 5 \end{pmatrix}\quad \bareme{1}\quad$\quad et \quad $\overrightarrow{B C} \begin{pmatrix} 0-5\\ 3-5 \end{pmatrix}=\overrightarrow{B c} \begin{pmatrix} -5\\ -2 \end{pmatrix}$\quad \bareme{1}\quad\\
        $\overrightarrow{B A} \cdot \widehat{B C}=-2 \times(-5)+5 x(-2)=10-10=0$\quad \bareme{1}\quad
    \end{corrBox}
    \item Que peut-on en conclure pour le triangle ABC?\\
    \begin{corrBox}
        $B A=B C=\sqrt{2^{2}+5^{2}}=\sqrt{29}$\\
        donc $A B C$ est un triangle rectangle et isocèle.\quad \bareme{1}\quad
    \end{corrBox}
\end{enumMet}

\eject
\exercice{ - équation cartésienne - ( 2 points )}

Déterminer dans le cas ci-dessous, une équation cartésienne de la droite D qui passe par le point $A$ et qui a pour vecteur normal le vecteur $\vec{n}$\\
$A(3 ; 4)$ et $\vec{n}\left(\begin{array}{c}-3 \\ 5\end{array}\right)$

\begin{corrBox}
    on a : $-3 x+5 y+c=0$\quad \bareme{1}\quad\\
$A(3 ; 4):-3 \times 3+5 \times 4+C=0 \Longleftrightarrow C=-11$\quad \bareme{0,5}\quad\\
Donc: $-3 x+5 y-11=0$\quad \bareme{0,5}\quad\\
\end{corrBox}

\exercice{ - dérivabilité - ( 4 points )}

Soit $f$ la fonction définie sur $R \operatorname{par} f(x)=-3 x^2+2x+5$\\
Montrer que la fonction f est dérivable pour $\mathrm{x}=3$ et calculer $f^{\prime}(3)$

\begin{corrBox}
$$
\begin{aligned} f(3+h)-f(3) &=-3(3+h)^{2}+2(3+h)+5-(-16) \\ &=-27-18 h-3 h^{2}+6+2 h+2 \\ &=-18 h-3 h^{2}+2 h \\ &=-3 h^{2}-16 h \quad \bareme{2} \\ \lim _{h \rightarrow 0} \frac{f(3+h)-f(3)}{h} &=\lim _{h \rightarrow 0}-3 h-16=-16 \quad \bareme{1} \\ \text { Donc: } f^{\prime}(-3) &=-16 \quad \bareme{1}\end{aligned}
$$
\end{corrBox}

\exercice{ - dérivabilité - ( 5 points )}
Soit f la fonction définie sur $\left] 2 ;+\infty\right[$ par $ f(x)=\dfrac{1}{2x+5}$.\\
Montrer que la fonction f est dérivable pour $x=1$ et calculer $f^{\prime}(1)$

\begin{corrBox}
    $f(1+h)=\dfrac{1}{2(1+h)+5}=\dfrac{1}{2 h+7}\quad \bareme{1}\quad$\quad \quad \quad $f(1)=\dfrac{1}{7}\quad \bareme{1}\quad$\\
    $f(1+h)-f(1)=\dfrac{7-(2 h+7)}{7(2 h+7)}=\dfrac{2 h}{14 h+49}\quad \bareme{1}\quad$\\
    $\lim \limits_{h \rightarrow 0} \dfrac{f(1+h)-f(a)}{h}=\lim \limits_{h \rightarrow 0} \dfrac{2}{14 h+49}=\dfrac{2}{49}\quad \bareme{1}\quad$\\
    Donc $f^{\prime}(1)=\dfrac{2}{49}\quad \bareme{1}\quad$
\end{corrBox}

\end{document}

